\documentclass[twocolumn]{article}
\usepackage[T1]{fontenc}
\usepackage{amsmath}
\usepackage{graphicx}

\title{Capacitor Charging and the Time Constant}
\author{Joseph Manno}

\begin{document}
\maketitle

\begin{abstract}
This document introduces the charging behaviour of a resistor-capacitor (RC) network and derives the capacitor voltage at one time constant.
\end{abstract}


\section{Introduction}
This document introduces the charging behaviour of a resistor-capacitor (RC) network. When a capacitor \textit{C} is charged through a series resistor \textit{R} from a constant supply $V_0$, the voltage capacitor voltage rises according to

\begin{equation}
    V(t) = V_0\left(1 - e^{-t/RC}\right),
\end{equation}

where the produce \textit{RC} has units of seconds and is called the time constant, $\tau = RC$.

\section{Voltage at the Time Constant}
Substituting $t = \tau$ into(1):

\begin{align*}
V(\tau) &= V_0 \left( 1 - e^{-\tau/RC} \right) \\
&= V_0 \left( 1 - e^{-RC/RC} \right) \quad \text{since } \tau = RC \\
&= V_0 \left( 1 - e^{-1} \right) \\
&\approx V_0 (1 - 0.368) \\
&\approx 0.632\, V_0.
\end{align*}

The exponent cancelling to -1 in (2) is why the 0.632 fraction is universal - $R$ and $C$ dropped out of the calculation before the final line, so this fraction holds for any $RC$ network. Substituting $V_0 = 5.0V$:

\begin{equation}
V(\tau)\approx 0.632 x 5.0V \approx 3.16V,
\end{equation}

matching the dashed line in Figure 1. 

\begin{figure}[t]
\centering
\includegraphics[width=\columnwidth]
{rc_charging.pdf}
\caption{Charging curve for
$R = 10\,\text{k}\Omega$, $C = 22\,\text{pF}$.}
\label{fig:rc}
\end{figure}

\end{document}
